
\thispagestyle{empty}

%-----Si desea usar el archivo Campos_Plantilla_UdeA_LaTex_IEEE_2022_ para llenar la información, seleccione o digite los datos en el documento de Word cuyo texto está coloreado, luego selecciones todo el contenido y ejecute la acción copiar; finalmente ejecutando pegar, reemplace los comandos de este archivo desde la línea 8 hasta la línea 83

%--------------------reemplase desde aqui------------------
%TÍTULO
 \newcommand{\mititulo}{Diseño y Desarrollo de un Actuador Dinámico de Alta Precisión con Control de Desplazamiento, Velocidad y Aceleración}%---- elimine "Título y/o subtítulo del trabajo de grado" e ingrese aquí el título del documento
 
%TÍTULO "CORTO" (que aparece en los encabezados de página)
  \newcommand{\smalltitulo}{Diseño y Desarrollo de un Actuador Dinámico de Alta Precisión}%---- Reemplaza con las primeras palabras de tu título o palabras significativas del mismo, seguido de puntos suspensivos. No permitas que título y número de página pasen a una segunda línea o más. Si no es tan extenso escriba el título completo en el campo anterior sin puntos suspensivos
   
  
 \begin{center}

    \includegraphics[scale=0.72]{images/escudo_udea.png}

    \vspace{2cm}
    \textbf{\mititulo} \\[2cm]%-----(en minúscula, y mayúscula cuando lo amerite: nombres propios, siglas, etc.). Título en negrita.
    
    %-----Ejemplo de autores (en minúsculas; no omitir ningún nombre o apellido; no abreviar ni dejar solo iniciales; marcar las tildes correspondientes).
    Juan Camilo Gallego Bedoya
    
    \vspace{2cm}
    Trabajo de Grado para optar al título de  Físico%-----Si deseas incluir el género o géneros de los integrantes, lo ingresas manualmente. Ejemplos: Bibliotecóloga, Bibliotecólogas, Bibliotecólogos.
    
    \vspace{2cm}
    Asesor:\\
    Johny Alexander Jaramillo Gallego, Doctorado en Física%-----(en minúscula; no omitir ningún nombre o apellido; no abreviar ni dejar solo iniciales; marcar las tildes correspondientes).

    Co-Asesor:\\
    Fabian Andres Castaño Usuga, Doctorado en Física


    \vspace{1cm}
    Universidad de Antioquia\\
    Facultad de Ciencias Exactas y Naturales\\
    Pregrado\\
    Medellín UdeA\\
    2026 %------Determina con tu asesor cuál es el año para registrar: año de entrega del documento o año en que se recibe el título académico. 

\end{center}